% !TEX root = chapter-standalone.tex

\section{From DPIs to DPs}
\linkvideo{spring2021-design:design:from-dpi-to-dp} % From DPI to DP

A DPI (\cref{def:DPI}) describes a relation between three sets:~\funsp, \ressp, \impsp.
If we are not interested in the implementations, but just in the relation between \funsp and \ressp, then we can describe a DPI more compactly as a ``DP''\@.


\begin{remark}
    \label{rem:DP-from-DPI}
    Given a DPI~$\tup{\funsp, \ressp, \impsp, \prov,\req}$ it is always possible to obtain from it the following DP
    \begin{equation}
        \label{eq:dpi_to_dp_1}
        \defmapperiod{
            \adp
        }{
            \funspop \Ptimes \ressp
        }{
            \toinPos
        }{
            \Bool
        }{
            \tup{\fun\Fop,\res}
        }{
            \exists \imp \setin \impsp \colon (\fun \posleqof\funsp \prov(\imp)) \booland (\req(\imp)\posleqof\ressp \res)
        }
        %            \adp\colon \funsp\op \Ctimes \ressp & \toinPos  \Bool \\
        %            \tup{\fun\Fop,\res}                   & \mapsto \exists \imp \setin \impsp \colon (\fun \posleqof\funsp \prov(\imp)) \booland (\req(\imp)\posleqof\ressp \res).
    \end{equation}
    Evaluating this DP is the same as asking whether the set
    \begin{equation}
        \label{eq:dpi_to_dp_2}
        \makeset{ \imp \setin \impsp \colon (\fun \posleqof\funsp \prov(\imp)) \booland (\req(\imp)\posleqof\ressp \res)}
    \end{equation}
    is empty or not.
\end{remark}
\vfill\pagebreak
\begin{example}
    Recall \cref{exa:dpi_elmotor}, with the catalogue of electric motors in \cref{tab:electric_motors_3}.

    \begin{table*}[h]
        \centering
        \adjustbox{max width=\textwidth}{%
            \begin{tabular}{c|c|c|c|c|c}
                \F{$\unit[\text{Torque}]{[kg\cdot cm]}$} & \colI{Motor ID}        & \R{\unit[Weight]{[g]}} & \R{\unit[Max Power]{[W]}}
                                                         & \R{\unit[Cost]{[USD]}} \\
                \hline
                \F{0.18}                                 & \colI{\textsf{1204}}   & \R{60.0}               & \R{2.34}                  & \R{19.95} \\
                \F{0.95}                                 & \colI{\textsf{1206}}   & \R{140.0}              & \R{3.00}                  & \R{19.95} \\
                \F{0.65}                                 & \colI{\textsf{1207}}   & \R{130.0}              & \R{2.07}                  & \R{12.95} \\
                \F{3.7}                                  & \colI{\textsf{2267}}   & \R{285.0}              & \R{4.76}                  & \R{16.95} \\
                \F{1.9}                                  & \colI{\textsf{2279}}   & \R{165.0}              & \R{5.40}                  & \R{164.95} \\
                \F{19.0}                                 & \colI{\textsf{1478}}   & \R{1,000}              & \R{8.96}                  & \R{49.95} \\
                \F{2.2}                                  & \colI{\textsf{2299}}   & \R{150.0}              & \R{5.90}                  & \R{59.95}
            \end{tabular}%
        }
        \caption{A simplified catalogue of motors.
        }
        \label{tab:electric_motors_3}
    \end{table*}
    \begin{table*}[h]
        \centering
        \adjustbox{max width=\textwidth}{%
            \begin{tabular}{c|c|c|c|c|c}
                \F{$\unit[\text{Torque}]{[kg\cdot cm]}$} & \R{\unit[Weight]{[g]}} & \R{\unit[Max Power]{[W]}}
                                                         & \R{\unit[Cost]{[USD]}} \\
                \hline
                \F{0.18}                                 & \R{60.0}               & \R{2.34}                  & \R{19.95} \\
                \F{0.95}                                 & \R{140.0}              & \R{3.00}                  & \R{19.95} \\
                \F{0.65}                                 & \R{130.0}              & \R{2.07}                  & \R{12.95} \\
                \F{3.7}                                  & \R{285.0}              & \R{4.76}                  & \R{16.95} \\
                \F{1.9}                                  & \R{165.0}              & \R{5.40}                  & \R{164.95} \\
                \F{19.0}                                 & \R{1,000}              & \R{8.96}                  & \R{49.95} \\
                \F{2.2}                                  & \R{150.0}              & \R{5.90}                  & \R{59.95}
            \end{tabular}%
        }
        \caption{Feasibility relations for the \SY{design problem} of motors.
        }
        \label{tab:electric_motors_dp}
    \end{table*}

    \begin{figure}[tbh]
        \centering
        \includesag{520_dp_electric_motor}
        \caption{Electric motor design problem.}
        \label{fig:dp_em_2}
    \end{figure}

    The catalogue induces a \SY{design problem}~$\adpan{\mathrm{EM}}$, where each feasibility relation between functionality and resources is reported in \cref{tab:electric_motors_dp}.
    with diagrammatic form as in \cref{fig:dp_em_2}.
    In particular, we can query the \SY{design problem} for combinations of \textF{functionalities} and \textR{resources}.
    For instance:
    \begin{equation}
        \adpan{\mathrm{EM}}\pars{\F{\unit[0.2]{kg\cdot cm}}, \R{\tup{\unit[50.0]{g}},\R{\unit[2.0]{W}},\R{\unit[15.0]{USD}}} }=\false,
    \end{equation}
    since no listed model can provide~$\F{\unit[0.2]{kg\cdot cm}}$ \textF{torque} by requiring the set of \textR{resources}~$\tup{\R{\unit[50.0]{g}},\R{\unit[2.0]{W}},\R{\unit[15.0]{USD}}}$ or less.

\end{example}


We have already seen in \cref{rem:DP-from-DPI} that we can obtain a DP from a DPI.
We can make this more formal and say that there exists a \SY{forgetful semifunctor} from \DPI to \DP.

\begin{definition}
    \label{def:dpitodpsemi}
    The \SY{forgetful semifunctor}~$\dpitodp\colon \DPI\fto \DP$ is given by:
    \begin{enumerate}
        \item Identity on the objects:~$\dpitodpob(\posA)=\posA$.
        \item Given~$\adp=\tup{\funsp,\ressp,\impsp,\prov,\req}$, the action on morphisms is given by
              \begin{equation}
                  \defmapperiod{\dpitodpmor(\adp)}
                  {\funspop \Ptimes \ressp}
                  {\toinPos}
                  {\Bool}
                  {\tup{\fun\opel, \res}}
                  {\exists \imp \setin \setI\colon (\fun \posleqof{\funsp}\prov(\imp))\booland (\req(\imp)\posleqof{\ressp} \res)}
              \end{equation}
    \end{enumerate}
\end{definition}

\begin{lemma}
    \cref{def:dpitodpsemi} indeed defines a \SY{semifunctor}.
\end{lemma}
\begin{proof}
    Consider
    \begin{equation}
        \begin{aligned}
            \adpa & =\tup{\funposA,\resposB,\impspn{1},\provn{1},\reqn{1}}, \\
            \adpb & =\tup{\funposB,\resposC,\impspn{2},\provn{2},\reqn{2}},
        \end{aligned}
    \end{equation}
    We need to show
    \begin{equation}
        \dpitodpmor(\adpa \dpithenDPI \adpb)=\dpitodpmor(\adpa) \dpthenDP \dpitodpmor(\adpb).
    \end{equation}
    Let's start with the left term.
    One has
    \begin{equation}
        \label{eq:dpitodp_a}
        \begin{aligned}
            \dpitodpmor(\adpa \dpithenDPI \adpb)(\funposAel\F{\opel},\resposCel) & =
            \exists \imp \setin \impsp\colon (\funposAel \posleqof{\posA}\provn{1}(\impn{1}))\booland (\reqn{2}(\impn{2})\posleqof{\posC} \resposCel),
        \end{aligned}
    \end{equation}
    where~$\impsp=\makeset{\impn{1}\tupconcat \impn{2}\setin \makecprod{\impspn{1}, \impspn{2}} \mid \reqn{1}(\impn{1})\posleqof{\posB}\provn{2}(\impn{2})}$.

    On the other hand,
    \begin{equation}
        \label{eq:dpitodp_b}
        \begin{aligned}
             & (\dpitodpmor(\adpa) \dpthenDP \dpitodpmor(\adpb))(\funposAel\F{\opel},\resposCel) \\
             & =\bigvee_{\posBel \setin \posBset}\dpitodpmor(\adpa)(\funposAel\F{\opel},\resposBel)\booland \dpitodpmor(\adpb)(\funposBel\F{\opel},\resposCel) \\
             & =\bigvee_{\posBel \setin \posBset} \pars{{\exists \impn{1} \setin \impspn{1}\colon (\funposAel \posleqof{\posA}\provn{1}(\impn{1}))\booland (\reqn{1}(\impn{1})\posleqof{\posB} \resposBel)}} \\
             & \booland \pars{{\exists \impn{2} \setin \impspn{2}\colon (\funposBel \posleqof{\posB}\provn{2}(\impn{2}))\booland (\reqn{2}(\impn{2})\posleqof{\posC} \resposCel)}}
        \end{aligned}
    \end{equation}
    Consider the following cases:
    \begin{itemize}
        \item If~$\dpitodpmor(\adpa \dpithenDPI \adpb)(\funposAel\F{\opel},\resposCel)=\true$, there exist~$\impn{1}\setin \impspn{1},\impn{2}\setin \impspn{2}$ for which
              \begin{equation}
                  \begin{aligned}
                      \funposAel         & \posleqof{\posA}\provn{1}(\impn{1}), \\
                      \reqn{2}(\impn{2}) & \posleqof{\posC} \resposCel, \\
                      \reqn{1}(\impn{1}) & \posleqof{\posB}\provn{2}(\impn{2}).
                  \end{aligned}
              \end{equation}
              The first two terms are clear, and the last term implies that there exists a~$\posBel\setin \posB$ such that
              \begin{equation}
                  (\funposBel\posleqof{\posB}\provn{2}(\impn{2}))
                  \booland (\reqn{1}(\impn{1})\posleqof{\posB}\resposBel),
              \end{equation}
              implying~$(\dpitodpmor(\adpa) \dpthenDP \dpitodpmor(\adpb))(\funposAel\F{\opel},\resposCel)=\true$.
        \item The case
              \begin{equation}
                  \prftree{\dpitodpmor(\adpa \dpithenDPI \adpb)(\funposAel\F{\opel},\resposCel)=\false}{(\dpitodpmor(\adpa) \dpthenDP \dpitodpmor(\adpb))(\funposAel\F{\opel},\resposCel)=\false}
              \end{equation}
              follows analogously.
        \item The other direction is easier to show, since clearly
              \begin{equation}
                  \prftree{(\dpitodpmor(\adpa) \dpthenDP \dpitodpmor(\adpb))(\funposAel\F{\opel},\resposCel)=\true}{\dpitodpmor(\adpa \dpithenDPI \adpb)(\funposAel\F{\opel},\resposCel)=\true}
              \end{equation}
              and
              \begin{equation}
                  \prftree{(\dpitodpmor(\adpa) \dpthenDP \dpitodpmor(\adpb))(\funposAel\F{\opel},\resposCel)=\false}{\dpitodpmor(\adpa \dpithenDPI \adpb)(\funposAel\F{\opel},\resposCel)=\false}
              \end{equation}
              by inspecting \cref{eq:dpitodp_a} and \cref{eq:dpitodp_b}.
    \end{itemize}
\end{proof}

In the other direction, we can take a DP and find a corresponding DPI.

Given a DP~$\adp\colon \funspop \Ptimes \ressp \toinPos \Bool$, we can define a DPI $\tup{\funsp,\ressp, \impsp,\prov,\req}$
by setting
\begin{equation}
    \begin{aligned}
        \impsp & =\makeset{\tup{\fun,\res}\setin \funsp \cartprod \ressp \colon \adp(\fun\F{\opel},\res)}, \\
        \prov  & \colon \tup{\fun,\res}\mapsto \fun, \\
        \req   & \colon \tup{\fun,\res}\mapsto \res,
    \end{aligned}
\end{equation}
However, this operation is not a semifunctor, since it does not preserve composition.

\devel {

    We obtain another \SY{semifunctor}.

    \begin{definition}
        \label{def:dptodpisemi}
        The \SY{semifunctor}~$\dptodpi\colon \DP\fto \DPI$ is given by:
        \begin{enumerate}
            \item Identity on the objects:~$\dpitodpob(\posA)=\posA$.
            \item Given~$\adp\colon \funspop \Ptimes \ressp \toinPos \Bool$, the action on morphisms is given by
                  \begin{equation}
                      \dptodpimor(\adp)=\tup{\funsp,\ressp, \impsp,\prov,\req},
                  \end{equation}
                  where
                  \begin{equation}
                      \begin{aligned}
                          \impsp & =\makeset{\tup{\fun,\res}\setin \funsp \cartprod \ressp \colon \adp(\fun\F{\opel},\res)}, \\
                          \prov  & \colon \tup{\fun,\res}\mapsto \fun, \\
                          \req   & \colon \tup{\fun,\res}\mapsto \res,
                      \end{aligned}
                  \end{equation}
        \end{enumerate}
    \end{definition}

    \begin{lemma}
        \cref{def:dptodpisemi} indeed defines a \SY{semifunctor}.
    \end{lemma}
    \todotext{
        @Gioele: the last passage in the proof is wrong. (the two sets have different dimensions!)
    }
    \begin{proof}
        Consider~$\adpa\colon \funposA\Ptimes \resposB \toinPos \Bool$ and~$\adpb\colon \funposB\Ptimes \resposC\toinPos \Bool$.
        We need to show
        \begin{equation}
            \dptodpimor(\adpa \dpthenDP \adpb)=\dptodpimor(\adpa)\dpithenDPI \dptodpimor(\adpb).
        \end{equation}
        Let's start with the left term.
        One has
        \begin{equation}
            \dptodpimor(\adpa \dpthenDP \adpb)=\tup{\funposA,\resposC,\impsp,\prov,\req},
        \end{equation}
        where
        \begin{equation}
            \begin{aligned}
                \impsp & =\makeset{\tup{\funposAel,\resposCel}\setin \funposA \cartprod \resposC \colon (\adpa \dpthenDP \adpb)(\funposAel,\resposCel)}, \\
                       & =\makeset{\tup{\funposAel,\resposCel}\setin \funposA \cartprod \resposC \colon \bigvee_{\posBel\setin \posB}\adpa(\funposAopel,\resposBel)\booland \adpb(\funposBopel,\resposCel)} \\
                \prov  & =\projstart, \\
                \req   & =\projend,
            \end{aligned}
        \end{equation}
        where~$\projstart$ is the projection mapping returning the first element of a tuple, and~$\projend$ is the projection mapping returning the last element of a tuple.
        For the right-hand side, instead, one has
        \begin{equation}
            \dptodpimor(\adpa) \dpithenDPI \dptodpimor(\adpb)=\tup{\funposA,\resposC,\impsp\colI{'},\prov\F{'},\req\R{'}},
        \end{equation}
        where
        \begin{equation}
            \begin{aligned}
                \impsp\colI{'} & =\makeset{\impn{1}\tupconcat \impn{2}\setin \makecprod{\impspn{1}, \impspn{2}}\mid \reqn{1}(\impn{1})\posleqof{\posB}\provn{2}(\impn{2})} \\
                \impspn{1}     & =\makeset{\tup{\funposAel,\resposBel}\setin \funposA \cartprod \resposB \colon \adpa(\funposAel,\resposBel)}, \\
                \impspn{2}     & =\makeset{\tup{\funposBel,\resposCel}\setin \funposB \cartprod \resposC \colon \adpb(\funposBel,\resposCel)}, \\
                \provn{1}      & =\projstart \\
                \provn{2}      & =\projstart \\
                \reqn{1}       & =\projend \\
                \reqn{2}       & =\projend \\
                \prov\F{'}     & =\projstart \\
                \req\R{'}      & =\projend.
            \end{aligned}
        \end{equation}
        One can already notice that~$\prov=\prov\F{'}$ and~$\req=\req\R{'}$.
        To finish the proof, we can massage~$\impsp\colI{'}$:
        \begin{equation}
            \begin{aligned}
                \impsp\colI{'} & =\makeset{\impn{1}\tupconcat \impn{2}\setin \makecprod{\impspn{1}, \impspn{2}}\mid \reqn{1}(\impn{1})\posleqof{\posB}\provn{2}(\impn{2})} \\
                               & =\makeset{\tup{\funposAel,\resposBel,\funposBel\F{'},\resposCel}\setin \funposA \cartprod \resposB\cartprod \funposB\cartprod \resposC \colon \adpa(\funposAel\F{\opel},\resposBel) \booland \adpb(\funposBel\F{\opel},\resposCel)
                \booland \resposBel \posleqof{\posB}\funposBel\F{'}} \\
                               & =\makeset{\tup{\funposAel,\resposCel}\setin \funposA \cartprod \resposC \colon \bigvee_{\posBel\setin \posB}\adpa(\funposAopel,\resposBel)\booland \adpb(\funposBopel,\resposCel)} \\
                               & =\impsp.
            \end{aligned}
        \end{equation}

    \end{proof}
}


